\cotefile{156-4}
\scanpage{42}\mspage{41}
\uline{Prop} Soit $F'$ un raffinement de $F$. Conditions équivalentes

a) $F'$ est une subdivision

b) Il n'existe pas de $X\in\unsure{\mathcal M|M}$ qui raffine $F$, et disjoint de $F'$.

c) \deleted{Il n'existe pas \unsure{Ou bien} $\widetilde F'=\widetilde F$, \unsure{sinon}, \illegible{1} \unsure{de} \unsure{il} il n'existe pas de lieu du $F$} Il n'existe pas de lieu dans l'ombre $\widetilde F$, qui soit disjoint de $F'$ i.e. \inserted{(\deleted{\unsure{par}}\unsure{Cond|Corol} 6)} $\forall\deleted{\unsure{x}}x\unsure{\in|+}\widetilde F\setminus\widetilde F'$, $\exists\,y\in\widetilde F'$ tel que $x$ et $y$ \uline{non compatibles}, ou (positivement) : Si $x\in\widetilde F$ est compatible à tout élément de $\widetilde F'$, il est dans $\widetilde F'$.

\uline{Dém} $a\Longrightarrow b$, car si $X\in\mathcal M$ raffine $F$ et est disjoint de $F'$, alors la figure $F\cup F_X$ est un raffinement de $F$, \unsure{strictmt|strictement} plus grand que $F'$,

$b)\Longrightarrow c)$ tautologique

$c)\Longrightarrow a)$ \unsure{Si l'on|Si on} \unsure{avait|en avait}, $F'\unsure{\subset|\in} \uline{F''}\unsure{\not\ll|\ll}F$, soit $X\in F''$, tel que $X\notin F'$, et soit $x\in\widetilde X^\circ$. Alors $x\in\widetilde F\setminus\widetilde F'$, et \deleted{je dis} que $x$ est compatible à $F'$ donc c) n'est pas vérifié.
% slanted
\margin{à vérifier}

Mais on a utilisé que $x$ est compatible à $F'$, il faut un axiome pour l'assurer.

\uline{Cond} 7 Soient $X,Y\in\mathcal M$, \inserted{$X$ et $Y$ compatibles} $X\ne Y$. \inserted{(donc $\widetilde X^\circ\cap\widetilde Y^\circ=\varnothing$)} Alors tout $x\in\widetilde X^\circ$ est disjoint de tout $y\in\widetilde Y^\circ$ (ou, ce qui revient au $=$ par Cond 6, \deleted{tout} $x\in X^\circ$ est disjoint de $Y$)
% Three vertical emphasis strokes beside the last three lines of Cond 7.

\uline{Cor. 1} \unsure{Soient} $F$ une figure, $F'$ une sous-figure et $X\in F\setminus F'$. Alors $\unsure{\forall|\mathrm{tt}}\,x\in X^\circ$ est \deleted{compatible} \inserted{disjoint} de $F'$.
\scanpage{43}\mspage{42}
\uline{Cor. 2} Soit $F$ une figure. Deux lieux dans \deleted{\uline{Cor. 2} Soit $F$ une figure. Deux} l'ombre de $F$, qui appartiennent à des ombres ouvertes de strates distinctes, sont \uline{disjoints}.

Soient $\deleted{X}F$ une figure, \deleted{$X\in\unsure{\mathcal M|M}$ un raffinement de $F$,} \inserted{$F'$ un raffinement} de $F$. Je \unsure{veux|vaux} \unsure{definir|définir} une application
\[
 F'\longrightarrow F
\]
en associant à \unsure{tt|tout} $X'\in F'$ \unsure{\{l'unique\}|(l'unique)} $X\in F$ tel que
\[
 \widetilde{X'}^\circ\subset\widetilde X^\circ
\]
(ou encore le \uline{plus petit} \deleted{\unsure{des|x|X}} $\unsure{X'|X}\in F$ tel que $\widetilde{X'}^\circ\subset\widetilde X$). \unsure{Savons-nous} déjà qu'un tel $X$ existe, i.e. que \unsure{la|R} relation d'équiv. canonique sur l'ombre de $F'$, est plus fine que celle induite par l'ombre de $F$ ? Soit $X$ un élément de $F$ \uline{minimal}, parmi ceux satisfaisant $X'\ll X$ (existe, à cause de l'axiome de dim), \unsure{savons nous|savons-nous} que dès lors $\widetilde{X'}^\circ\subset\widetilde X^\circ$ ? \deleted{\unsure{Soit}} On sait que $\widetilde{X'}\subset\widetilde X$ donc si $x\in X'^\circ\inserted{\subset\widetilde{X'}}$, on a $x\in\widetilde{\deleted{\unsure{X|Y}}Y}^\circ\subset\widetilde Y$ pour $Y\unsure{\leq|\subseteq}X$, on veut prouver que $Y=X$. Il semblerait qu'il manque quelque chose.
\scanpage{44}\mspage{43}
\uline{Cond 8} Soit $X,Y,X'$ dans $\unsure{\mathcal M|M}$, $Y\ll X$, $X'<X$. Considérons l'ens des $Y'\in\deleted{\mathcal M}F_Y$ (i.e. $Y'\leq Y$) \unsure{tels|telles} que $Y'\ll\unsure{X|x}$ :
% slanted
\margin{\unsure{Axiome} du raffinement induit.}
\diagram{\begin{tikzcd}[ampersand replacement=\&,row sep=0.1em,column sep=0.35em]
Y \& \ll \& X\\
\tikz[baseline=-.5ex]\node[rotate=90,inner sep=0pt] {$\leq$}; \& {} \& \tikz[baseline=-.5ex]\node[rotate=90,inner sep=0pt] {$\unsure{<|\leq}$};\\
Y' \& \ll \& X'
\end{tikzcd}}

(lesquels forment une sous-figure $R_{X'}$ de \unsure{la|le} multistrate $\deleted{\unsure{F}}F_Y$, \deleted{qu'on appellera} raffinant $X'$, et qu'on appellera le raffinement de $X'$, induit par $Y$ sur $X'$). Alors \unsure{tout\deleted{e}|toute} raffinement commun $Z$ de $Y$ et de $X'$ raffine un des $Y'$, i.e. raffine $R_{X'}$. (Donc Toute figure qui raffine à la fois $\deleted{\unsure{Y|F}}F_Y$, et $F_{X'}$, raffine $R_{X'}$, qui est \deleted{\unsure{donc} la borne} borne inf. de $F_Y$ et $F_{X'}$ pour la relation $\ll$) En particulier,
\[
 :\quad \widetilde Y\cap\widetilde{X'}=\bigcup_{Y'\in R_{X'}}\widetilde{Y'}
\]

\uline{Cor} \deleted{\unsure{Soit} $x\in\widetilde{\unsure{Y|X}}^\circ$,} Supposons que $\deleted{\unsure{\exists}}\nexists\,X_1\in\mathcal M$, avec $Y\ll X_1<X$ (c'est la situation de \unsure{tantôt}) Alors $\widetilde Y^\circ\subset\widetilde X^\circ$.

En effet, si $x\in\widetilde Y^\circ$, \unsure{si on|s'on} avait $x\in\widetilde{X'}$ pour $X'<X$, on aurait $x\in\widetilde{Y'}$ avec $Y'\in R_{X'}$, donc $Y'\unsure{=|\Rightarrow}Y$ par déf. de $\widetilde Y^\circ$, donc $Y\ll X'_{\deleted{\unsure{1}}}<X$, \inserted{contraire à \unsure{hyp.}}
\scanpage{45}\mspage{44}
\uline{Corollaire} Si $\deleted{\unsure{X}}Y\ll X$, alors \deleted{la relation} toute ombre ouverte d'une \deleted{\unsure{strate de}} \inserted{strate \unsure{eff.|aff.}} $Y$ est contenue \unsure{dans|ds} l'ombre ouverte d'une strate de $X$. Plus généralement et par suite, pour \unsure{tte} figure $F$ et \deleted{\unsure{t}} \deleted{sous-figure} \inserted{raffinement} $F'$ de $F$, \unsure{toute} \deleted{strate} \deleted{\unsure{de} ombre ouverte d'une} strate ouverte de $\unsure{F'|F}$ est contenue dans une (unique) strate ouverte de $F$.

\deleted{\uline{\unsure{Cond}}. \unsure{8} \unsure{Si} \unsure{/}}\enspace$\backslash$

Je viens de m'apercevoir que dans le cas du pseudo disque, \deleted{\unsure{une}} deux lieux disjoints \deleted{\unsure{cq}} \inserted{$x,y$} de l'intérieur (lesquels définissent donc deux ``raffinements'' dudit pseudo-disque), qu'il existe une ``Subdivision'' du disque dont ils soient des \unsure{noeuds|nœuds}. [\uline{NB}. Les \uline{noeuds} d'une figure \inserted{$\deleted{\unsure{F}}F$} sont les multistrates \inserted{de $\deleted{\unsure{X|F}}\unsure{F}$} réduites à un lieu, i.e. les éléments de $F\cap\unsure{\mathcal L|L}$] Donc \deleted{\unsure{la no}} il n'y a pas lieu de considérer $\{x\unsure{,|.}y\}$ comme un ``raffinement'', vu qu'il ne s'étend pas en une subdivision. Donc il est inapproprié
\scanpage{46}\mspage{45}
de définir une figure $F'$ comme \unsure{un} \uline{\unsure{raffinant|raffinement}} de $F$, si chaque $X'\in F'$ \deleted{\unsure{déf}} est un raffinement d'un $X\in F$. Donc il vaut mieux prendre comme notion \unsure{primitive|première} celle de \uline{subdivision} d'une \deleted{\unsure{mul}} figure, (qui \unsure{va|en} se \unsure{ramener|ramène} à celle d'une multistrate), et définir un raffinement d'une figure comme sous-figure d'une subdivision.

Donc, nouveau départ !

Une \unsure{contrée} $C$ est définie par un triple
\[
 (\unsure{\mathcal M|M},\unsure{\mathcal F|\mathscr F},\Sigma)
\]
où $\mathcal M$ est un \unsure{ens} (ens. des \uline{multistrates}), $\mathcal F$ un ensemble de parties de $\mathcal M$ ($\mathcal F\unsure{\in|\subset}\unsure{\mathscr P|\mathfrak P}(\mathscr P(\mathcal M))$), \inserted{\unsure{appelées|appelés} \uline{figures},} et $\Sigma$ une relation sur $\mathcal F$ dite : $F'$ \uline{subdivise} $F$. \unsure{Axiomes|Axiome}
% slanted
\margin{Si $F\in\mathcal F$, les $X\in F$ sont les (multi)strates de $F$}

\uline{C 1} Pour tout $X\in\mathcal M$, $\exists$ \unsure{une|un} plus petit\unsure{e|$\subset$} \deleted{partie} figure \deleted{\unsure{qui}} $F_X\unsure{\subset|{,}}$ telle que $X\in F$.

On l'appelle \unsure{la} \uline{figure élémentaire} définie par $X$. On écrit
\[
 X\leq Y\ \overset{\mathrm{def}}{\Longleftrightarrow}\ F_X\subset F_Y\qquad\text{i.e. }X\in F_Y
\]
et on trouve une relation d'\uline{ordre}$^{*}$ $\leq$ sur $\mathcal M$.\textsuperscript{$*$}
% Two vertical strokes run beside C 1.
% slanted
\margin{$^{*}$ On veut que ce soit une relation d'ordre et pas seulement de préordre, i.e. \uline{C$'$1} si \deleted{\unsure{pr}} $X\ne Y$, $\exists$ $F\in\mathcal F$ qui \unsure{contient} un des deux \unsure{ms} $X,Y$, et pas l'autre, i.e. $F_X\ne F_Y$.}
% slanted
\margin{\unsure{--} i.e. $\mathcal F$ sépare les éléments de $\mathcal M$}
\scanpage{47}\mspage{46}
% slanted
\margin{Ainsi, $\unsure{\mathcal M|M}$ s'identifie à un ens de parties de $\unsure{\mathcal F|\mathscr F}$ (savoir les $F_X$) i.e. à l'ens des figures élémentaires. \deleted{Donc on peut considérer que la structure est \unsure{portée} par $\mathcal F$, muni de $\mathcal F_{\unsure{\mathrm{él}|\mathrm{el}}}$ et de $\unsure{\Sigma}$, la relation de subdivision (relation d'ordre)} \deleted{Il y a \unsure{une sorte de symétrie|deux sortes de symétries} -- on travaille sur $\mathcal F$, \unsure{ou bien|comme} \inserted{par \unsure{les axiomes des}} \unsure{et tantôt} sur $\mathcal M$,  1°) relation d'ordre $\unsure{C'|\leq}$ (subdivision, ou $\ll$ raffinement) et ens. de parties (élémentaires, les figures).}}
\deleted{\unsure{Cela étant}}

\deleted{$\unsure{X}\ll\unsure{Y}\quad\unsure{F'}\unsure{\supset|\subset}F_X$} Si $F$ est une figure, alors $X\in F_{\deleted{\unsure{X|x}}}$ équivaut à $F_X\subset F$, qui implique que pour tout $Y\in F_X$ \inserted{i.e. $Y\leq X$} on a $F_Y\subset F$ i.e. $Y\in F$. Ainsi, $F$ est une partie de $\mathcal M$ \uline{fermée} pour $\leq$.

On dit que deux figures \inserted{$F,G$} \uline{sont compatibles} si $F\cup G$ est une figure, et que $X,Y\in\mathcal M$ sont \uline{compatibles} si $F_X,F_Y$ le sont

\uline{C \deleted{\unsure{2|3}}\unsure{3|2}} Pour que $F$ et $G$ soient compatibles, il f. et s. que $\forall\,\unsure{x|X}\in F,\ Y\in G$, $X$ et $Y$ le \unsure{soient|sont}.
% Two vertical emphasis strokes.
% slanted
\margin{correspond à Fig \unsure{4|9}}

\uline{Cor} \deleted{En particulier si $X$} Pour que $X,Y\in\mathcal M$ soient compatibles, il f et s. qu'ils \uuline{app}artiennent à une figure

\deleted{\uline{C 3} Si $F$ est.}
% slanted
\margin{\tikz[baseline=(a.base)]{\node[inner sep=1pt] (a) {\unsure{intervertir}}; \draw[<->] (a.north east) .. controls +(-.7,.2) and +(-.7,-.2) .. (a.south east);}}

\uline{C \deleted{\unsure{3|2}}\unsure{2|3}} Soit $F$ une figure, alors tte partie de $F$ fermée pour $\leq$ est aussi une figure (i.e. pour qu'une partie $F'$ de $F$ soit une figure, il f et il \uline{suffit} qu'elle soit fermée pour $\leq$).
% Two vertical emphasis strokes.

(\uline{Sous-figures}) Si $\mathcal M\ne\varnothing$ cela \deleted{\unsure{signifie}} \unsure{implique} que la partie vide de $\mathcal M$ est une figure. Si $\mathcal M=\varnothing$, \deleted{cela implique que} \inserted{on aura $\mathcal F=\varnothing$} \deleted{si $\mathcal F$} ou $\mathcal F=\{\varnothing\}$, et on suppose que c'est le deuxième cas : \uline{figure vide} existe tjrs
% slanted
\margin{\uline{Contrée vide} si $\mathcal M=\varnothing$}
\scanpage{48}\mspage{47}
\uline{Définition} Soient $F,\deleted{\unsure{F'}}G$ \unsure{deux|des} figures. On dit que $\deleted{\unsure{F'}}G$ est un \uline{raffinement} de $F$, s'il existe une \deleted{\unsure{raffinée|raffin}} subdivision $F'$ de $F$, telle que $G\subset F'$ \inserted{sous-figure}.
% slanted
\margin{On écrit $F\ll G$}
% Irregular ink blots in the blank space below the marginal formula.

On veut que ce soit une relation d'ordre, et pour ceci on va poser

\uline{C 4} Soient $F$ une figure, $G$ une sous-figure \inserted{(i.e. $G$ figure et $G\unsure{\subset|\in}F$)}, $G'$ une subdivision de $G$. Alors \deleted{il existe une subdivision $F'$ de $F$,} \inserted{$G'\cup(F\setminus G)$ est une} subdivision \inserted{de $F$}, et la réunion est disjointe (i.e. \deleted{\unsure{hors de $G$}} % the two strokes beside C 4 extend to the boxed formula
\[
 \boxed{G'\cap F\subset G}
\]
\deleted{donc} $G'\cap F=G'\cap G$).
% slanted
\margin{C'est \unsure{démenti|évidemment} \unsure{par exemple|faux pour} les \unsure{orbites|strates} fermées et strates \unsure{augmentées}\ldots\ \uline{cf plus} bas.}
% slanted
\margin{i.e \unsure{toute} multistrate d'une subdivision de $G$ qui est \uline{compatible} avec $F$ \uline{est dans $G$}. OK}
\margin{OK $\longrightarrow$}
\margin{Corriger : $\exists$ subdivision qui induit $G'$ $+$ la relation $G'\cap\unsure{F\to\subset|F'\subset|F\subset}G$}

\uline{Cor. 1} \uline{Toute} \deleted{raffinement} \inserted{\uline{subdivision}} \uline{d'une sous figure de $F$ est induite par une subdivision de $F$}.
% A long curved stroke groups the two lines.

\uline{Cor 2} La relation de raffinement entre figures
\[
 F\ll G
\]
est une relation d'ordre.
% slanted
\margin{\unsure{OK}.}

\deleted{\uline{\unsure{C 5}} \inserted{\uline{Proposition}} Soient $X,\deleted{\unsure{Y}}\in\unsure{\mathcal M|M}$, \inserted{$\deleted{\unsure{F}}G\in\unsure{\mathcal F|\mathscr F}$.} Pour que $X\unsure{\in|\notin}\deleted{\unsure{F_Y}}G$ \deleted{i.e. $F_X\subset\unsure{F_Y|F_X}$} il f. et s. que $X$ et $\deleted{\unsure{Y}}G$ compatibles (i.e. $F=F_X\cup\deleted{\unsure{F_Y}}G$ une figure) et $X$ \deleted{raffine} $\inserted{\ll}\deleted{\unsure{F_Y}}G$ i.e $F_X$ raffine $\deleted{\unsure{F_Y}}G$.}
% slanted
\margin{C'est la \deleted{\unsure{Condition}} \inserted{\unsure{définition}} de $\leq$ via $\ll$ et \unsure{R}}

\deleted{\deleted{\unsure{On}} Supposons en effet que $X$ et $\unsure{F|G}$ compatibles \deleted{il \unsure{faut} $X\unsure{,|\in}Y\in F$}, $F$ figure, et que $F_X$ raffine la sous-figure $F_Y\deleted{\unsure{=}} \deleted{\unsure{G.}}$. Donc $\exists$ subd. $G'$ de}
\scanpage{49}\mspage{48}
\deleted{$G$ telle que $F_{X'}\subset G'$ i.e. $X'\in G'$. Cette subdivision est induite par subdivision $G'\cup(F\setminus G)$, qui est disjointe, et comme}
\[
\deleted{X'\in G'}
\]
\deleted{\unsure{Supposons} \unsure{un|une}}

\deleted{\uline{Proposition}}\inserted{Cor. 3} Soient $\deleted{\unsure{F|G}}G,F$ deux figures. Pour que $\deleted{F}G\subset F$\inserted{ (i.e. $G$ sous figure de $F$)}, il f et s. que $\deleted{F}G$ \uline{compatible} à $F$ (i.e. $\deleted{\unsure{G|F}}F\cup\deleted{F}G$ une figure, soit $\deleted{\unsure{F|G}}H$) et $\deleted{F}G\ll F$, i.e $G\subset\deleted{G'}F'$, $\deleted{\unsure{G|F}}F'$ sub. de $F$.

En effet, par \deleted{\unsure{Cond|C}} C 4, \deleted{comme}
\diagram{\begin{tikzcd}[ampersand replacement=\&,row sep=.15em,column sep=.25em]
F \& \subset \& H=F\cup G \\
F'\arrow[u,Rightarrow,"\text{subd.}"] \& {} \& {}\\
\tikz[baseline=-.5ex]\node[rotate=90,inner sep=0pt] {$\subset$}; \& {} \& {}\\
\deleted{\unsure{F|G}}G \& {} \& {}
\end{tikzcd}}
% The following calculation is written immediately to the right of the diagram.
\[
G\subset \underset{\deleted{\unsure{F'\cap H|G'\cap H}}}{\deleted{G'}F'\cap H}\ \deleted{\unsure{\subset G'\cap H|\subset F'\cap H}}\ \underset{\mathrm C\,4}{\mathrel{\tikz[baseline=-.35ex]\draw[line cap=round] (.35,.16) .. controls (.12,.18) and (.02,.08) .. (0,0) -- (1.18,0);}}\ F
\]
donc $G\subset F$.

En particulier, on trouve :
% slanted
\margin{C'est la partie \unsure{sérieuse} de C 4 \(\left(\vphantom{\dfrac{1}{1}}\right.\)}
\[
\text{Si } X\in\unsure{\mathcal M|M},\ F\in\unsure{\mathcal F|\mathscr F},\ \text{alors}\quad X\in F\ \Longleftrightarrow\
\left\{\begin{aligned}
&X\ \text{comp. à }F\\
&X\ \text{raffine }F
\end{aligned}\right.
\]
donc
\[
\text{Si } X,Y\in\mathcal M,\ \text{alors}\quad X\leq Y\ \Longleftrightarrow\
\left\{\begin{aligned}
&X\ \text{comp. à }Y\\
&X\ \deleted{\text{raffine}}\ll Y
\end{aligned}\right.
\]

% slanted
\margin{\uline{NB} $F$ et $G$ disj. \unsure{ssi} $X\in F,\ Y\in G\Rightarrow X$ et $Y$ disj.$^{\mathrm{ts}}$ (\unsure{i.e.|ou} ou non compatibles, \unsure{ou} $X$ et $Y$ \unsure{elt|et} \uuline{distincts})}

\deleted{\uline{\unsure{C 5|C 2}}} $F,G\in\mathcal F$ sont dites \uline{disjointes} si elles sont compatibles (i.e $F\cup G\in\mathcal F$) et $F\cap G=\varnothing$. On dit que $X,Y$ sont disjointes, \unsure{ssi} $F_X$ et $F_Y$ le sont i.e. $X$ et $Y$ compatibles, et $\nexists Z$ avec $Z\leq X,Y$.

Si $F$ et $G$ sont disjointes, alors pour $F'\subset F,G'\subset G$ sous-figures, elles sont disjointes. De plus ;
